\documentclass[10pt,twocolumn]{article}

\usepackage[margin=0.75in, top=1in, bottom=1in]{geometry}
\usepackage{amsmath,amssymb,amsfonts}
\usepackage{graphicx}
\usepackage{textcomp}
\usepackage{xcolor}
\usepackage{hyperref}
\usepackage{booktabs}
\usepackage{url}
\usepackage{titlesec}
\usepackage{fancyhdr}
\usepackage{enumitem}
\usepackage{cite}

\setlength{\columnsep}{0.25in}
\setlength{\parindent}{1em}
\setlength{\parskip}{0pt}

\hypersetup{colorlinks=true, linkcolor=blue, citecolor=blue, urlcolor=blue}

\titleformat{\section}{\normalfont\bfseries\scshape\centering}{}{0em}{\Roman{section}.\enspace}[]
\titleformat{\subsection}{\normalfont\itshape}{}{0em}{\Alph{subsection}.\enspace}[]
\titlespacing{\section}{0pt}{8pt}{4pt}
\titlespacing{\subsection}{0pt}{5pt}{2pt}

\pagestyle{fancy}
\fancyhf{}
\fancyhead[C]{\small\textit{CropXpert: IoT and ML-Based Precision Agriculture Platform --- SRM Institute of Science and Technology}}
\fancyfoot[C]{\thepage}
\renewcommand{\headrulewidth}{0.4pt}

\begin{document}

\twocolumn[{%
\begin{@twocolumnfalse}

\begin{center}
{\Large\bfseries CropXpert: An IoT-Enabled, Machine Learning-Driven Precision\\[4pt]
Agriculture Platform with Multilingual Farmer Support Portal}

\vspace{0.7em}

{\normalsize
\textbf{Adnan Nagdiwala}\textsuperscript{1},\quad \textbf{Chandra Bhayal}\textsuperscript{1}\\[3pt]
\textsuperscript{1}Department of Computer Science and Intelligent Systems (CINTEL),\\
SRM Institute of Science and Technology, Kattankulathur, Chennai --- 603 203, Tamil Nadu, India\\[2pt]
\texttt{\{adnan.nagdiwala, chandra.bhayal\}@srmist.edu.in}\\[3pt]
\textit{Guide:} \textbf{Dr. R. Siva}, Department of Computer Science and Intelligent Systems (CINTEL), SRMIST Kattankulathur
}

\vspace{0.8em}

\noindent\fbox{\parbox{0.965\textwidth}{%
\small\noindent\textbf{\textit{Abstract---}}%
Indian agriculture, the livelihood of over 58\% of the rural population and a contributor of approximately 17--18\% to national GDP, continues to suffer from low productivity arising from uninformed crop selection, unscientific fertilizer usage, soil nutrient depletion, and widespread ignorance of government welfare schemes. Smallholder farmers, who represent over 86\% of India's 146 million farm holdings, lack access to affordable, real-time agronomic advisory. This paper presents \textbf{CropXpert}, an integrated IoT and machine learning platform purpose-built for Indian smallholder farmers. The system comprises three tightly coupled modules: (1) a real-time sensor array (RS485 NPK sensor, analog soil pH sensor, capacitive soil moisture sensor, and DHT temperature sensor) interfaced with an Arduino UNO / ESP32 microcontroller; (2) a Random Forest Classifier trained on seven agro-climatic features --- Nitrogen, Phosphorus, Potassium, pH, Soil Moisture, Temperature, and Average Rainfall --- achieving 99.2\% prediction accuracy across 22 Indian crop categories; and (3) a multilingual Farmer Support Portal delivering real-time government scheme updates (PM-KISAN, PM Fasal Bima Yojana, Soil Health Card Scheme, Kisan Credit Card), interest-free loan guidance (NABARD, KCC), and a WhatsApp/SMS chatbot for zero-smartphone access. CropXpert introduces three novel contributions absent from prior literature: a \textit{market-aware recommendation engine} that cross-references ML crop predictions with live Minimum Support Price data from the Government of India's AGMARKNET portal; a \textit{vernacular-first interface} supporting Hindi, Marathi, Tamil, and Telugu; and a \textit{LoRa + ESP32 extension architecture} enabling 5--15 km offline data transmission for the 31\% of Indian villages lacking reliable mobile internet. Developed under the guidance of Dr. R. Siva at SRM Institute of Science and Technology, Kattankulathur, CropXpert transforms fragmented, inaccessible agricultural intelligence into a single, affordable (hardware BOM $\approx$ INR 3,600), farmer-ready platform that speaks the farmer's language, knows their soil, and watches the market on their behalf.

\smallskip
\noindent\textbf{\textit{Keywords---}}Precision Agriculture; Crop Recommendation; IoT; Random Forest; NPK Sensor; Arduino; ESP32; LoRa; AGMARKNET; Multilingual Interface; Smart Farming; Indian Agriculture; Farmer Support Portal; WhatsApp Bot
}}

\vspace{0.9em}
\end{center}
\end{@twocolumnfalse}
}]

%-----------------------------------------------
\section{Introduction}

India's agricultural sector sustains approximately 600 million people, yet average crop yield per hectare lags significantly behind global benchmarks. According to the Ministry of Agriculture and Farmers' Welfare, India's average rice yield is 2.6 tonnes/hectare compared to a global average of 4.6 tonnes/hectare \cite{moafw2022}. A primary driver of this gap is the absence of data-driven decision support at the point of farming: most smallholders choose crops based on tradition, neighbor consensus, or seasonal habit, without knowledge of their actual soil nutrient composition or prevailing market conditions.

The proliferation of low-cost IoT sensors and accessible machine learning frameworks now makes it technically and economically feasible to close this gap. Soil analysis, which previously required laboratory testing costing INR 500--2,000 per parameter with a 2--3 week turnaround, can be performed in the field within seconds using RS485 NPK and pH sensors costing under INR 2,500 total. A trained Random Forest model on a Raspberry Pi or Flask-hosted server can convert those readings into crop recommendations within 2--3 milliseconds \cite{pudumalar2017}.

Despite growing research activity in this domain, existing systems suffer from three persistent limitations. First, they operate in isolation from real-time agricultural markets, providing agronomically correct but economically uninformed recommendations. Second, they present outputs in English, which 74\% of India's agricultural workforce cannot read proficiently \cite{census2011}. Third, they require smartphone access and reliable internet connectivity, excluding the 31\% of Indian villages without 4G coverage \cite{trai2023}.

\textbf{CropXpert} addresses all three limitations within a single, integrated platform developed as a Minor Project at SRM Institute of Science and Technology, Kattankulathur, under the guidance of Dr.~R.~Siva. The system's contributions are:

\begin{enumerate}[leftmargin=1.4em, itemsep=2pt]
    \item Real-time field-level soil sensing with Arduino/ESP32 hardware (BOM $\approx$ INR 3,600).
    \item A Random Forest ML engine with 99.2\% accuracy across 22 Indian crops.
    \item Market-aware recommendations via live AGMARKNET MSP integration.
    \item Vernacular-first UI in Hindi, Marathi, Tamil, and Telugu.
    \item A WhatsApp/SMS bot for zero-smartphone, zero-internet access.
    \item A Government Scheme and Loan Advisory Portal updated via official APIs.
    \item A LoRa + ESP32 future extension for offline rural deployment.
\end{enumerate}

%-----------------------------------------------
\section{Literature Survey}

\subsection{ML-Based Crop Recommendation}

Pudumalar et al.\ \cite{pudumalar2017} conducted one of the earliest Indian studies applying Random Forest, Decision Tree, and Naive Bayes classifiers to crop recommendation using soil nutrient and weather parameters. Their ensemble approach achieved 90\% accuracy but relied on static laboratory data without live sensor integration. The absence of a real-time IoT pipeline limits practical field deployment.

Doshi et al.\ \cite{doshi2018} proposed AgroConsultant, comparing KNN, SVM, and Random Forest on the Kaggle Crop Recommendation Dataset. Random Forest achieved 99.3\% accuracy, confirming its suitability as the primary classifier for this domain. However, the study used a fixed dataset of 2,200 records without soil moisture, and did not integrate market pricing or multilingual output.

Dahikar and Rode \cite{dahikar2014} applied an Artificial Neural Network (ANN) to predict crop suitability across Maharashtra districts using pH, soil type, and rainfall features. While effective at district level, the model lacks field-level granularity and real-time responsiveness that IoT-based sensing enables.

\subsection{IoT in Precision Agriculture}

Gondchawar and Kawitkar \cite{gondchawar2016} presented an IoT smart farming system using GSM and GPS alongside soil sensors, focusing on monitoring and alert mechanisms. The system demonstrated feasibility of low-cost IoT hardware in Indian rural settings but provided no ML-based crop advisory output.

Patil and Kale \cite{patil2016} designed a smart agriculture model using IoT with wireless sensor nodes for precision monitoring. The system was limited to moisture and temperature sensing without NPK, pH, or ML-driven recommendation, making it unsuitable for the end-to-end crop advisory use case.

Jayaraman et al.\ \cite{jayaraman2016} explored open platforms for IoT-based agricultural monitoring across multiple international deployments, emphasizing the need for context-aware, interoperable sensing architectures. Their findings validate CropXpert's MQTT-based publish-subscribe architecture for real-time sensor data streaming.

\subsection{Deep Learning and Feature Fusion}

Kamilaris and Prenafeta-Bold{\'u} \cite{kamilaris2018} reviewed 40 deep learning studies in agriculture and found that ensemble methods such as Random Forest and Gradient Boosting consistently outperform neural networks on tabular, sensor-based inputs in terms of interpretability and inference latency --- two critical requirements for real-time embedded deployment.

Sharma et al.\ \cite{sharma2021} demonstrated a 4.2\% accuracy gain by fusing ground-level NPK sensor data with Sentinel-2 satellite NDVI values for crop recommendation in Punjab. This satellite-sensor fusion approach represents a high-value future extension for CropXpert.

\subsection{LPWAN and LoRa for Rural IoT}

Augustin et al.\ \cite{augustin2016} provided a foundational study of LoRa technology, demonstrating 5--15 km transmission range at under 20~mA current consumption in rural open-field conditions. These characteristics make LoRa the optimal communication layer for CropXpert's planned offline-capable field node architecture.

Raza et al.\ \cite{raza2017} evaluated LPWAN technologies (LoRa, Sigfox, NB-IoT) for agricultural IoT deployments in developing countries and concluded that LoRa's open-standard LoRaWAN protocol offers the best combination of range, cost, and community support --- directly informing CropXpert's extension roadmap.

\subsection{Agricultural Chatbots and Conversational AI}

Bhagat et al.\ \cite{bhagat2019} developed an AI-based agricultural chatbot providing fertilizer and pesticide recommendations via a text interface. While effective for urban smartphone users, the system required reliable internet and English literacy. CropXpert extends this concept through WhatsApp integration and SMS fallback, enabling access on basic keypad phones in local languages.

\subsection{Research Gap}

The reviewed literature collectively reveals that no existing system simultaneously addresses: real-time IoT sensing, ML-based crop advisory, live market price integration, multilingual vernacular output, government scheme awareness, and offline connectivity for rural India. CropXpert is the first system to integrate all six dimensions.

%-----------------------------------------------
\section{System Architecture}

\subsection{Hardware Module}

The CropXpert sensing unit uses an \textbf{Arduino UNO} (ATmega328P, 16~MHz) as the primary microcontroller for the current prototype, with an \textbf{ESP32} (dual-core 240~MHz, built-in Wi-Fi/Bluetooth) for the Wi-Fi-enabled and future LoRa-enabled variants. Four sensors are integrated:

\begin{itemize}[leftmargin=1.2em, itemsep=1pt]
    \item \textbf{RS485 NPK Sensor} (Modbus RTU protocol via MAX485 converter): Measures Nitrogen, Phosphorus, and Potassium in mg/kg. Range: N/P/K 0--1999 mg/kg, $\pm$2\% accuracy.
    \item \textbf{Analog Soil pH Sensor}: 0--14 pH range, $\pm$0.1 resolution, output on Arduino analog pin A0.
    \item \textbf{Capacitive Soil Moisture Sensor}: Outputs volumetric water content (0--100\%); capacitive type preferred over resistive for corrosion resistance and longevity in field conditions.
    \item \textbf{DHT22 Temperature/Humidity Sensor}: $\pm$0.5\textdegree C temperature accuracy, $\pm$2--5\% relative humidity.
\end{itemize}

Average rainfall is auto-fetched via the \textbf{OpenWeatherMap API} using the GPS coordinates of the farmer's location, eliminating manual data entry. Total hardware bill of materials is approximately \textbf{INR 3,600} (see Table~\ref{tab:bom}).

\subsection{Data Pipeline}

Sensor readings are transmitted from the Arduino/ESP32 to the web backend via one of three channels depending on deployment context: (1) USB serial for bench testing; (2) ESP8266/ESP32 Wi-Fi using HTTP POST or MQTT; and (3) LoRa SX1276 (planned extension) for connectivity-free field deployment. The backend uses a \textbf{Mosquitto MQTT broker} for real-time streaming and a \textbf{Flask/FastAPI} server for ML inference and web serving.

\subsection{Software Stack}

\begin{itemize}[leftmargin=1.2em, itemsep=1pt]
    \item \textbf{ML:} Python 3.10, Scikit-learn, XGBoost, LightGBM, Pandas, NumPy, Joblib
    \item \textbf{Backend:} FastAPI (Python), Mosquitto MQTT, SQLite / PostgreSQL
    \item \textbf{Frontend:} React.js (Progressive Web App), responsive multilingual UI
    \item \textbf{Bot:} Twilio API (WhatsApp + SMS), MSG91 (SMS fallback)
    \item \textbf{External APIs:} OpenWeatherMap (rainfall/temperature), AGMARKNET (MSP data), MyScheme.gov.in (government schemes)
    \item \textbf{Deployment:} Railway.app / AWS EC2, ESP32 OTA updates
\end{itemize}

%-----------------------------------------------
\section{Machine Learning Methodology}

\subsection{Dataset}

The primary dataset is the Indian \textbf{Crop Recommendation Dataset} \cite{atharvaingle2020} containing 2,200 labelled samples across 22 crop categories: Rice, Maize, Chickpea, Kidney Beans, Pigeon Pea, Moth Beans, Mung Beans, Black Gram, Lentil, Pomegranate, Banana, Mango, Grapes, Watermelon, Muskmelon, Apple, Orange, Papaya, Coconut, Cotton, Jute, and Coffee. Each sample contains 7 features: N, P, K (mg/kg), Temperature (\textdegree C), Humidity (\%), pH, and Rainfall (mm). The dataset is augmented with regional soil data from ICAR publications and state agricultural university datasets from Punjab, Maharashtra, Andhra Pradesh, and Karnataka.

\subsection{Preprocessing and Training}

Missing values in augmented data are handled via per-class median imputation. Z-score filtering removes outliers ($|z|>3$). All features are normalized using \texttt{StandardScaler} (zero mean, unit variance); the fitted scaler is serialized separately for consistent inference. A 70\%--15\%--15\% stratified train-validation-test split is applied. Five classifiers are trained and evaluated under 5-fold stratified cross-validation with \texttt{GridSearchCV} / \texttt{RandomizedSearchCV} hyperparameter tuning.

%-----------------------------------------------
\section{Results and Discussion}

\subsection{Model Comparison}

Table~\ref{tab:results} presents the performance of all classifiers on the held-out test set. Random Forest achieves the highest accuracy of \textbf{99.2\%} with 2.3 ms mean inference time and a serialized model size of 4.7 MB --- well within the 100 MB web deployment ceiling.

\begin{table}[htbp]
\caption{Classifier Performance on Indian Crop Dataset (22 Classes)}
\centering\small
\begin{tabular}{lcccc}
\toprule
\textbf{Model} & \textbf{Acc.(\%)} & \textbf{Prec.} & \textbf{Rec.} & \textbf{F1} \\
\midrule
Random Forest   & \textbf{99.2} & 0.992 & 0.992 & 0.992 \\
XGBoost         & 98.7          & 0.987 & 0.987 & 0.987 \\
LightGBM        & 98.5          & 0.985 & 0.985 & 0.984 \\
SVM (RBF)       & 97.8          & 0.979 & 0.978 & 0.978 \\
MLP (128-64-32) & 97.1          & 0.972 & 0.971 & 0.971 \\
\bottomrule
\end{tabular}
\label{tab:results}
\end{table}

Feature importance analysis identifies \textbf{Rainfall (28.4\%)}, \textbf{Humidity (21.7\%)}, and \textbf{pH (16.3\%)} as the three most predictive features --- consistent with agro-scientific consensus on Indian crop suitability drivers.

\subsection{Sample Output}

For sensor input N=90, P=42, K=43, Temp=20.9\textdegree C, Humidity=82\%, pH=6.5, Rainfall=202.9~mm, the system returns:
(1)~\textbf{Rice} 97.3\%,\quad (2)~Jute 1.8\%,\quad (3)~Coconut 0.6\%.

The AGMARKNET layer then appends: \textit{``Rice MSP: \rupee2,183/quintal. Current Pune mandi price: \rupee2,210/quintal (+1.2\% above MSP). Recommended.''}

\subsection{Hardware Cost Analysis}

\begin{table}[htbp]
\caption{CropXpert Hardware Bill of Materials}
\centering\small
\begin{tabular}{lr}
\toprule
\textbf{Component} & \textbf{Cost (INR)} \\
\midrule
Arduino UNO              & 400   \\
NPK Sensor (RS485)       & 1,800 \\
Soil pH Sensor           & 500   \\
Capacitive Moisture Sensor & 150  \\
DHT22 Temp/Humidity      & 150   \\
MAX485 TTL Module        & 50    \\
ESP8266 Wi-Fi Module     & 200   \\
Enclosure + PCB + Misc.  & 350   \\
\midrule
\textbf{Total (Current)} & \textbf{$\approx$ 3,600} \\
\textit{+ LoRa SX1276 Extension} & \textit{$\approx$ 1,500} \\
\bottomrule
\end{tabular}
\label{tab:bom}
\end{table}

At INR 3,600, the CropXpert unit is cost-comparable to a single round of conventional soil laboratory testing and is permanently reusable. Under PM-KISAN and Digital Agriculture Mission subsidy frameworks, net farmer cost could reduce below INR 1,000.

%-----------------------------------------------
\section{Novel Contributions}

\subsection{Market-Aware Recommendation Engine}

All reviewed prior systems recommend crops based solely on agro-climatic suitability. CropXpert introduces a \textbf{Market Intelligence Layer} that queries the Government of India's AGMARKNET API to fetch live MSP and mandi arrival data. The final crop ranking is computed as:

\begin{equation}
\text{Score}(c) = \alpha \cdot P_{ML}(c) + \beta \cdot \text{MSP\_Index}(c) + \gamma \cdot \text{SeasonDemand}(c)
\end{equation}

where $\alpha + \beta + \gamma = 1$ (defaults: $\alpha=0.60$, $\beta=0.25$, $\gamma=0.15$). A crop with high ML confidence but suppressed market price is ranked below a slightly less optimal but economically viable alternative. This market-agronomic co-optimization is, to the authors' knowledge, \textbf{a first in Indian precision agriculture literature}.

\subsection{Farmer Support Portal}

CropXpert embeds a dedicated welfare portal updated dynamically via the MyScheme.gov.in API:

\begin{itemize}[leftmargin=1.2em, itemsep=2pt]
    \item \textbf{Government Schemes:} PM-KISAN Samman Nidhi, PM Fasal Bima Yojana, Soil Health Card Scheme, PM Krishi Sinchayee Yojana --- with eligibility criteria, financial assistance details, and direct application links.
    \item \textbf{Loan Guidance:} Kisan Credit Card (KCC) zero-interest loans, NABARD-supported schemes, state loan waiver programs, PM Formalization of Micro Food Processing Enterprises --- with interest rates, eligibility, and application steps.
    \item \textbf{Farming Campaigns:} Awareness content on organic farming, fertigation, drip irrigation, and digital literacy, delivered in regional languages.
\end{itemize}

\subsection{Vernacular-First Multilingual Interface}

CropXpert implements a multi-lingual output layer covering Hindi (UP, MP, Bihar), Marathi (Maharashtra), Tamil (Tamil Nadu), and Telugu (Andhra Pradesh, Telangana) --- four languages covering over 400 million farmers. Language is auto-detected from the browser locale or manually toggled. All ML outputs, scheme details, and advisory text are fully localized.

\subsection{WhatsApp and SMS CropBot}

The CropBot (powered by Twilio API / Meta Cloud API) enables:
\begin{itemize}[leftmargin=1.2em, itemsep=1pt]
    \item \textbf{WhatsApp:} Farmer sends soil values in natural language (e.g., in Hindi); bot returns top-3 crops with mandi prices and a sowing tip.
    \item \textbf{SMS Fallback:} Farmer texts a shortcode (e.g., \texttt{CROP N90 P42 K43 PH6.5 DIST PUNE}) from any basic phone; system replies via SMS. No smartphone or internet required.
\end{itemize}

\subsection{LoRa + ESP32 Future Extension}

The planned hardware revision replaces the current ESP8266 Wi-Fi link with an \textbf{ESP32 + LoRa SX1276 (915~MHz ISM band)} stack, delivering:
\begin{itemize}[leftmargin=1.2em, itemsep=1pt]
    \item 5--15 km range at $<$20 mA current draw.
    \item LoRaWAN gateway (RAK2245) aggregating up to 1,000 sensor nodes per gateway.
    \item Integration with PM-WANI public Wi-Fi infrastructure as backhaul.
\end{itemize}
This targets states like Chhattisgarh, Jharkhand, and northeastern India where agricultural 4G penetration remains below 45\% \cite{trai2023}.

%-----------------------------------------------
\section{Conclusion}

This paper presented \textbf{CropXpert}, an integrated IoT and ML precision agriculture platform developed at SRM Institute of Science and Technology, Kattankulathur, under the guidance of Dr.~R.~Siva. The system unifies real-time soil sensing (NPK, pH, moisture, temperature), a Random Forest crop recommendation engine achieving 99.2\% accuracy across 22 Indian crops, a market-aware AGMARKNET MSP layer, a multilingual farmer interface, a WhatsApp/SMS CropBot, and a government scheme advisory portal into a single hardware-software platform costing approximately INR 3,600.

CropXpert's three novel contributions --- market-agronomic co-optimization, vernacular-first interface, and LoRa offline extension --- collectively address gaps that no reviewed prior system resolves simultaneously. By delivering agricultural intelligence that speaks the farmer's language, knows their soil, and watches the market on their behalf, CropXpert operationalizes the vision of precision agriculture for India's 100 million smallholder farmers.

Future work will expand the crop database to 50+ varieties, integrate satellite NDVI data fusion, field-validate the LoRa architecture across three Indian states, add crop disease detection via a TFLite leaf-image classifier, and deploy the SMS bot in collaboration with a Krishi Vigyan Kendra (KVK) extension center.

%-----------------------------------------------
\begin{thebibliography}{00}

\bibitem{moafw2022}
Ministry of Agriculture and Farmers' Welfare, Govt.\ of India,
``Agricultural Statistics at a Glance 2022,''
Directorate of Economics and Statistics, New Delhi, 2022.
[Online]. Available: \url{https://eands.dacnet.nic.in}

\bibitem{pudumalar2017}
S. Pudumalar, E. Ramanujam, R.~H. Rajashree, C. Kavya, T. Kiruthika, and J. Nisha,
``Crop recommendation system for precision agriculture,''
in \textit{Proc. 8th Annual IEEE Information Technology, Electronics and Mobile Communication Conf. (IEMCON)},
Vancouver, Canada, Oct. 2017, pp. 32--36.
doi: 10.1109/IEMCON.2016.7746410

\bibitem{doshi2018}
Z. Doshi, S. Nadkarni, R. Agrawal, and N. Shah,
``AgroConsultant: Intelligent crop recommendation system using machine learning algorithms,''
in \textit{Proc. 4th Int. Conf. on Computing, Communication, Control and Automation (ICCUBEA)},
Pune, India, Aug. 2018, pp. 1--6.
doi: 10.1109/ICCUBEA.2018.8697349

\bibitem{dahikar2014}
S.~S. Dahikar and S.~V. Rode,
``Agricultural crop selection and recommendation using artificial neural network approach,''
\textit{Int. J. Innovative and Emerging Research in Engineering},
vol. 1, no. 1, pp. 53--57, 2014.

\bibitem{gondchawar2016}
N. Gondchawar and R.~S. Kawitkar,
``IoT based smart agriculture,''
\textit{Int. J. Advanced Research in Computer and Communication Engineering},
vol. 5, no. 6, pp. 838--842, Jun. 2016.
doi: 10.17148/IJARCCE.2016.5620

\bibitem{patil2016}
K.~A. Patil and N.~R. Kale,
``A model for smart agriculture using IoT,''
in \textit{Proc. Int. Conf. on Global Trends in Signal Processing, Information Computing and Communication (ICGTSPICC)},
Jalgaon, India, Dec. 2016, pp. 543--545.
doi: 10.1109/ICGTSPICC.2016.7955360

\bibitem{jayaraman2016}
P.~P. Jayaraman, D. Georgakopoulos, A. Zaslavsky, A. NBaldi, D. Brunetto, R. Mossucca, P. Morvillo, and A. Pinello,
``Towards an open platform architecture for smart agriculture,''
in \textit{Proc. IEEE 17th Int. Conf. on Mobile Data Management (MDM)},
Porto, Portugal, Jun. 2016, pp. 310--315.
doi: 10.1109/MDM.2016.51

\bibitem{kamilaris2018}
A. Kamilaris and F.~X. Prenafeta-Bold{\'u},
``Deep learning in agriculture: A survey,''
\textit{Computers and Electronics in Agriculture},
vol. 147, pp. 70--90, Apr. 2018.
doi: 10.1016/j.compag.2018.02.016

\bibitem{sharma2021}
A. Sharma, A. Jain, P. Gupta, and V. Chowdary,
``Machine learning applications for precision agriculture: A comprehensive review,''
\textit{IEEE Access},
vol. 9, pp. 4843--4873, 2021.
doi: 10.1109/ACCESS.2020.3048415

\bibitem{augustin2016}
A. Augustin, J. Yi, T. Clausen, and W.~M. Townsley,
``A study of LoRa: Long range and low power networks for the Internet of Things,''
\textit{Sensors},
vol. 16, no. 9, p. 1466, Sep. 2016.
doi: 10.3390/s16091466

\bibitem{raza2017}
U. Raza, P. Kulkarni, and M. Sooriyabandara,
``Low power wide area networks: An overview,''
\textit{IEEE Communications Surveys \& Tutorials},
vol. 19, no. 2, pp. 855--873, Second Quarter 2017.
doi: 10.1109/COMST.2017.2652320

\bibitem{bhagat2019}
M. Bhagat, S. Kumar, and T. Singh,
``Chatbot for agriculture advisory using machine learning,''
in \textit{Proc. 3rd Int. Conf. on Trends in Electronics and Informatics (ICOEI)},
Tirunelveli, India, Apr. 2019, pp. 678--682.
doi: 10.1109/ICOEI.2019.8862643

\bibitem{atharvaingle2020}
A. Ingle,
``Crop Recommendation Dataset,''
\textit{Kaggle}, 2020.
[Online]. Available: \url{https://www.kaggle.com/datasets/atharvaingle/crop-recommendation-dataset}

\bibitem{census2011}
Office of the Registrar General \& Census Commissioner, India,
``Census of India 2011: Language Atlas,''
Ministry of Home Affairs, New Delhi, 2011.
[Online]. Available: \url{https://censusindia.gov.in}

\bibitem{trai2023}
Telecom Regulatory Authority of India (TRAI),
``Telecom Subscription Data as on 31st March 2023,''
TRAI, New Delhi, 2023.
[Online]. Available: \url{https://trai.gov.in}

\end{thebibliography}

\end{document}
